%
% chapter.tex -- Integration elementarer Funktionen
%
% (c) 2025 Prof Dr Andreas Müller
%
% !TeX spellcheck = de_CH
\chapter{Integration in geschlossener Form
\label{chapter:intalg}}
\kopflinks{Integration in geschlossener Form}
Im vorangegangenen Kapitel waren erfolgreich darin, eine Stammfunktion
einer rationale Funktionen zu finden.
Die Frage nach einer Stammfunktion von Funktionen wie den
Wurzelfunktionen, den trigonometrischen Funktionen oder einre Funktion
wie $e^{-x^2}$ bleibt aber weiterhin offen.
Der Hauptsatz der Infinitesimalrechnung garantiert zwar, dass es
eine reelle Stammfunktion gibt.
Dies scheint aber nicht die Antwort zu sein, die die Fragestellung
suggeriert.
Vielmehr wird wohl eine Stammfunktion als algebraischer Ausdruck
erwartet, der aus Polynomen, Wurzeln, Logarithmen, Exponentialfunktionen
und trigonometrischen Funktionen aufgebaut wird.
Um diese Aufgabe zu lösen ist erst zu klären, welche Funktionen
dazu überhaupt in Betracht zu ziehen sind und wie sie charakterisiert
werden können.
Der Begriff der \emph{elementaren Funktionen} wird in
Abschnitt~\ref{buch:intalg:section:elementar} eingeführt und es
wird gezeigt, dass sie mit den Exponential- und Logarithmusfunktionen
automatisch auch alle bekannten trigonometrischen und hyperbolischen
Funktionen umfassen.

Die angedeutete Aufgabe kann mit den elementaren Funktionen neu
formuliert werden: Welche elementaren Funktionen haben ein elementare
Stammfunktion?
Die Antwort gibt das \emph{Kriterium von Liouville}, welches in
Abschnitt~\ref{buch:intalg:section:liouville} 
vorgestellt wird.
Aus dem Kriterium lässt sich der \emph{Algorithmus von Risch}
ableiten, der in Abschnitt~\ref{buch:intalg:section:risch}
skizziert wird.
Zusammen mit heuristischen Methoden, Stammfunktionen zu finden,
bildet dieser Algorithmus die Basis der Integrationsfunktion
vieler Computeralgebrasysteme.

\input{chapters/070-intalg/1-elementar.tex}
%
% 2-wurzel.tex
%
% (c) 2026 Prof Dr Andreas Müller
%
\section{Integrale von $R(x,\vartheta)$ mit $\vartheta=\!\sqrt{ax^2+bx+c}$
\label{buch:intalg:section:wurzel}}
\kopfrechts{Integrale von $R(x,\vartheta)$ mit $\vartheta=\!\sqrt{ax^2+bx+c}$}
In Kapitel~\ref{chapter:ratint} ist es uns bereits gelungen, Stammfunktionen
für beliebige rationale Funktionen zu finden.
Der Formalismus von Abschnitt~\ref{buch:intalg:section:elementar} zeigt,
dass die Stammfunktion immer ein elementare Funktion ist.
In diesem Abschnitt soll untersucht werden, was sich ändert, wenn man den
Körper $\mathbb{Q}(z)$ der rationalen Funktionen algebraisch um die
Quadratwurzel 
\[
\vartheta = \!\sqrt{ax^2+bx+c}
\]
erweitert.
Die Vorgehensweise soll später verallgemeinert werden und zeigen, wie
allgemein entschieden werden kann, ob ein Integrand eine elementare
Stammfunktion hat.

%
% Der Funktionekörper
%
\subsection{Der Funktionenkörper $\mathbb{Q}(z,\vartheta)$
\label{buch:intalg:wurzel:subsection:koerper}}
Der Körper $\mathbb{Q}(z,\vartheta)$ besteht aus den Funktionen
\[
f(z)
=
\frac{
\tilde{p}_0 + \tilde{p}_1\vartheta + \tilde{p}_2\vartheta^2
+ \ldots + \tilde{p}_n\vartheta^n
}{
\tilde{q}_0 + \tilde{q}_1\vartheta + \tilde{q}_2\vartheta^2
+ \ldots + \tilde{q}_m\vartheta^m
}
\]
mit $p_k\in\mathbb{Q}(z)$, $k=1,\dots,n$, und $q_k\in\mathbb{Q}(z)$,
$k=1,\dots,m$.
Da $\vartheta$ eine algebraisches Element über $\mathbb{Q}$ ist, welches
die Identität
\[
\vartheta^2 = az^2 +bz +c
\]
mit $a,b,c\in\mathbb{Q}$ erfüllt, können die höheren Potenzen
$\vartheta^{2l}$ und $\vartheta^{2l+1}$ durch Potenzen des Radikanden
ersetze werden:
\begin{align*}
\vartheta^{2l}
&=
(az^2 + bz + c)^l
&&\text{und}&
\vartheta^{2l+1}
&=
(az^2 + bz + c)^l \vartheta.
\end{align*}
Es ist daher möglich, die Funktion $f(z)$ in der einfacheren Form
\begin{equation}
f(z)
=
\frac{p_0+p_1\vartheta}{q_0+q_1\vartheta}
\label{buch:algint:wurzel:eqn:funktion}
\end{equation}
zu schreiben.
Gesucht ist also eine Stammfunktion von
\[
f
\in
\mathbb{Q}(z,\vartheta)
=
\biggl\{
\frac{p_0+p_1\vartheta}{q_0+q_1\vartheta}
\;
\bigg|
\;
p_k,q_k\in\mathbb{Q}(z)\, k=0,1
\biggr\}
\]
als elementare Funktion.
Da schon die Stammfunktionen von rationalen Funktionen nicht unbedingt
rational sein müssen, kann man nicht erwarten, dass eine Stammfunktion
einer Funktion in $\mathbb{Q}(z,\vartheta)$ in diesem Körper liegen wird.
Es werden also Erweiterungen des Körpers um weitere Funktionen nötig sein,
um das Problem lösen zu können.

%
% Rationalisierung
%
\subsection{Rationalisierung
\label{buch:intalg:wurzel:subsection:rationalisierung}}
Die Darstellung~\eqref{buch:algint:wurzel:eqn:funktion} der Funktion kann
weiter vereinfacht werden, mit potenzieller Vereinfachung der Berechnung
der Stammfunktion.
Dazu wird \eqref{buch:algint:wurzel:eqn:funktion} mit
$q_0-q_1\vartheta$ zu
\begin{align}
f
&=
\frac{p_0+p_1\vartheta}{q_0+q_1\vartheta}
\cdot
\frac{q_0-q_1\vartheta}{q_0-q_1\vartheta}
\notag
\\
&=
\frac{
p_0q_0 +(p_0q_1+p_1q_0)\vartheta + p_1q_1\vartheta^2
}{
q_0^2-q_1^2\vartheta^2
}
\notag
\\
&=
\frac{
p_0q_0+p_1q_1(az^2+bz+c)
}{
q_0^2-q_1^2(az^2+bz+c)
}
+
\frac{
p_0q_1+p_1q_0
}{
q_0^2-q_1^2(az^2+bz+c)
}
\vartheta
\notag
\\
&=
p + q\vartheta
\label{buch:intalg:wurzel:eqn:linkombtheta}
\end{align}
erweitert.
Daraus liest man ab, dass der Funktionenkörper sogar als
\[
\mathbb{Q}(z,\vartheta)
=
\{
p+q\vartheta
\mid
p,q\in\mathbb{Q}(z)
\}
\]
geschrieben werden kann.

Die Relation $\vartheta^2-(az^2+bz+c)=0$ für das algebraische Element
$\vartheta$ bedeutet auch, dass
\[
\vartheta
=
\frac{az^2+bz+c}{\vartheta}
\]
und damit, dass der Funktionenkörper auch aus Linearkombinationen
\begin{equation}
\mathbb{Q}(z,\vartheta)
=
\biggl\{
p+q\frac{1}{\vartheta}
\bigg|
p,q\in\mathbb{Q}
\biggr\}
\label{buch:intalg:wurzel:eqn:linkombthetainverse}
\end{equation}
von $1$ und $1/\vartheta$ mit Koeffizienten in $\mathbb{Q}(z)$ erzeugen.

%
% Das Integrationsproblem
%
\subsection{Das Integrationsproblem
\label{buch:intalg:wurzel:subsection:integration}}
Die in Abschnitt~\ref{buch:intalg:wurzel:subsection:rationalisierung}
erfolgte Rationalisierung zeigt, dass eine rationale Funktion der Form
$R(z,\vartheta)$ immer in die Form \eqref{buch:intalg:wurzel:eqn:linkombtheta}
oder \eqref{buch:intalg:wurzel:eqn:linkombthetainverse} gebracht
werden kann.
Es müssen also nur noch Integrale der Form
\begin{equation*}
\int
p(z)+q(z)\vartheta
\,dz
\qquad\text{oder}\qquad
\int
p(z)+q(z)\frac{1}{\vartheta}
\,dz
\end{equation*}
gelöst werden, wobei $p$ und $q$ rationale Funktionen in $\mathbb{Q}(z)$
sind.
Das Integral über den ersten Summanden ist das Integral einer rationalen
Funktion und kann daher mit der Methode von Kapitel~\ref{chapter:ratint}
integriert werden.
Es bleibt daher nur noch Integrale der Form
\begin{equation}
\int q(z)\,\vartheta\,dz
\qquad\text{oder}\qquad
\int \frac{q(z)}{\vartheta}\,dz
\label{buch:intalg:wurzel:eqn:integrand}
\end{equation}
zu bestimmen.


Im Folgenden bevorzugen wir die zweite Form des Integrals in
\eqref{buch:intalg:wurzel:eqn:integrand}.
Um den Grund zu verstehen, betrachten wir die Ableitung
\begin{align*}
\vartheta^2
&=
az^2+bz+c
\\
2\vartheta\vartheta'
&=
2az + b
\intertext{und drücken $\vartheta'$ durch $\vartheta$ als}
\vartheta'
&=
\frac{2ax+b}{2\vartheta}
=
\frac{ax+b/2}{\vartheta}
&&\Rightarrow&
\frac{1}{\vartheta}
&=
\frac{\vartheta'}{ax+b/2}
\end{align*}
aus.
Man kann also $1/\vartheta$ durch die Ableitung $\vartheta'$ ersetzen,
was die Möglichkeit der partiellen Integration eröffnet.

Die rationale Funktion $q(z)$ im rechten Integral von
\eqref{buch:intalg:wurzel:eqn:integrand} ist ein Bruch von Polynomen
von $z$.
Durch Dividieren mit Rest lässt es sich schreiben als Summe
\begin{equation}
q(z)
=
\varphi(z)
+
w(z)
\qquad
\varphi(z)\in\mathbb{Q}[z],
w(z)\in\mathbb{Q}(z)
\label{buch:intalg:wurzel:eqn:summe}
\end{equation}
eines Polynoms und einer vollständig gekürzten rationalen Funktion,
deren Nenner grösseren Grad hat als der Zähler.

%
% Integranden der Form Polynom/theta
%
\subsection{Integranden der Form $\text{Polynom}/\vartheta$}
Der Summand $\varphi$ in
\eqref{buch:intalg:wurzel:eqn:summe}
lässt sich auf ein Integral nur von $1/\vartheta$ reduzieren,
wie der folgende Satz zeigt.

\begin{satz}
Sei $\varphi\in\mathbb{Q}[z]$ ein Polynom mit rationalen Koeffizienten
und $\vartheta=\!\sqrt{az^2+bz+c}$.
Dann gibt es ein Polynom $\psi\in\mathbb{Q}[z]$ mit 
\begin{equation}
\int\frac{\varphi}{\vartheta}
=
\psi\vartheta
+
A\int\frac{1}{\vartheta},
\label{buch:intalg:wurzel:eqn:polytheta}
\end{equation}
wobei $A$ eine Konstante ist.
\end{satz}

\begin{proof}
Wenn es das Polynom $\psi$ und die Konstante $A$ gibt, dann 
muss auch die Ableitung der Gleichung~\eqref{buch:intalg:wurzel:eqn:polytheta} 
erfüllt sein, also
\begin{equation}
\frac{\varphi}{\vartheta}
=
\psi'\vartheta + \psi\vartheta'+A\frac{1}{\vartheta}.
\label{buch:intalg:wurzel:eqn:polytheta'}
\end{equation}
Die Ableitung von $\vartheta$ ist
\[
\vartheta'
=
\frac{2az+b}{2\vartheta}
=
\frac{az+b/2}{\vartheta}.
\]
Eingesetzt in 
\eqref{buch:intalg:wurzel:eqn:polytheta'} und mit $\vartheta$ multipliziert
erhält man die Gleichung
\begin{align*}
\varphi
&=
\psi'\vartheta^2 + \psi (az+b/2) + A
\\
\varphi
&=
\psi' (az^2+bz+c) + \psi (az+b/2) + A
\end{align*}
Um diese Differentialgleichung für die Funktion $\psi$ zu lösen, verwenden
wir Ansätze der Form
\[
\renewcommand{\arraycolsep}{1pt}
\begin{array}{rcrcrcrcrcccrcr}
\varphi&=& \varphi_mz^m &+& \varphi_{m-1}z^{m-1} &+&   \varphi_{m-2}z^{m-2} &+&   \varphi_{m-3}z^{m-3} &+& \ldots &+& \varphi_1 z &+& \varphi_0\phantom{.}\\
\psi   &=&              & &    \psi_{m-1}z^{m-1} &+&      \psi_{m-2}z^{m-2} &+&      \psi_{m-3}z^{m-3} &+& \ldots &+&    \psi_1 z &+&    \psi_1\phantom{.}\\
\psi'  &=&              & &                      & & (m-1)\psi_{m-1}z^{m-2} &+& (m-2)\psi_{m-2}z^{m-3} &+& \ldots &+&   2\psi_2 z &+&    \psi_1.
\end{array}
\]
Durch Ausmultiplizieren und Koeffizientenvergleich entstehen die Gleichungen
\[
\renewcommand{\arraycolsep}{1pt}
\def\i#1{\rlap{$\scriptstyle #1$}\phantom{m-2}}
\begin{array}{lcrcrcrcrcrcrcr}
\varphi_{m  } &=& (m-1)a\psi_{m-1}   & &                    & &                    &+& a\psi_{m-1}   & &                   & &   \\
\varphi_{m-1} &=& (m-2)a\psi_{m-2}   &+& (m-1)b\psi_{m-1}   & &                    &+& a\psi_{m-2}   &+& (b/2)\psi_{m-1} & &   \\
\varphi_{m-2} &=& (m-3)a\psi_{m-3}   &+& (m-2)b\psi_{m-2}   &+& (m-1)c\psi_{m-1}   &+& a\psi_{m-3}   &+& (b/2)\psi_{m-2}   & &   \\
\varphi_{m-3} &=& (m-4)a\psi_{m-4}   &+& (m-3)b\psi_{m-3}   &+& (m-2)c\psi_{m-2}   &+& a\psi_{m-4}   &+& (b/2)\psi_{m-3}   & &   \\
              &\vdots &              & &                    & &                    & &               & &                   & &   \\
\varphi_{3}   &=&     2a\psi_{\i{2}} &+&     3b\psi_{\i{3}} &+&     4c\psi_{\i{4}} &+& a\psi_{\i{2}} &+& (b/2)\psi_{\i{3}} & &   \\
\varphi_{2}   &=&      a\psi_{\i{1}} &+&     2b\psi_{\i{2}} &+&     3c\psi_{\i{3}} &+& a\psi_{\i{1}} &+& (b/2)\psi_{\i{2}} & &   \\
\varphi_{1}   &=&                    & &      b\psi_{\i{1}} &+&     2c\psi_{\i{2}} &+& a\psi_{\i{0}} &+& (b/2)\psi_{\i{1}} & &   \\
\varphi_{0}   &=&                    & &                    & &      c\psi_{\i{1}} & &               &+& (b/2)\psi_{\i{0}} &+& A \\
\end{array}
\]
oder zusammengefasst
\[
\renewcommand{\arraycolsep}{1pt}
\def\i#1{\rlap{$\scriptstyle #1$}\phantom{m-2}}
\begin{array}{lcrcrcrcrcr}
\varphi_{m  } &=& \phantom{(}m\phantom{\mathstrut-0)}a\psi_{m-1}
                                     & &                             & &                    & &   \\
\varphi_{m-1} &=& (m-1)a\psi_{m-2}   &+&  ((2(m-1)+1)b/2)\psi_{m-1}  & &                    & &   \\
\varphi_{m-2} &=& (m-2)a\psi_{m-3}   &+&  ((2(m-2)+1)b/2)\psi_{m-2}  &+& (m-1)c\psi_{m-1}   & &   \\
\varphi_{m-3} &=& (m-3)a\psi_{m-4}   &+&  ((2(m-3)+1)b/2)\psi_{m-3}  &+& (m-2)c\psi_{m-2}   & &   \\
              &\vdots &              & &                             & &                    & &   \\
\varphi_{3}   &=&     3a\psi_{\i{2}} &+& 7b/2\phantom{)}\psi_{\i{3}} &+&     4c\psi_{\i{4}} & &   \\
\varphi_{2}   &=&     2a\psi_{\i{1}} &+& 5b/2\phantom{)}\psi_{\i{2}} &+&     3c\psi_{\i{3}} & &   \\
\varphi_{1}   &=&      a\psi_{\i{0}} &+& 3b/2\phantom{)}\psi_{\i{1}} &+&     2c\psi_{\i{2}} & &   \\
\varphi_{0}   &=&                    & &  b/2\phantom{)}\psi_{\i{0}} &+&      c\psi_{\i{1}} &+& A.\\
\end{array}
\]
Daraus kann man Rekursionsformeln für die Koeffizienten $\psi_k$ ableiten:
\begin{align*}
\psi_{m-1}
&=
\frac{1}{ma}\varphi_m
\\
\psi_{m-2}
&=
\frac{1}{(m-1)a}\biggl(
\varphi_{m-1}
-
((2(m-1)+1)b/2)\psi_{m-1}
\biggr)
\\
\psi_{m-3}
&=
\frac{1}{(m-2)a}\biggl(
\varphi_{m-2}
-
((2(m-1)+1)b/2)\psi_{m-2}
-
(m-1)c\psi_{m-1}
\biggr)
\\
&\phantom{a}\vdots
\\
\psi_{k}
&=
\frac{1}{(k+1)a}
\biggl(
\varphi_{k+1}
-
((2(k+1)+1)b/2)\psi_{k+1}
-
(k+2)
\psi_{k+2}
\biggr)
\\
&\phantom{a}\vdots
\\
\psi_0
&=
\frac{1}{a}\biggl(
\varphi_1
-
3b/2\psi_1
-
2c\psi_2
\biggr)
\\
A
&=
\varphi_0 -(b/2)\psi_0 - c\psi_1.
\end{align*}
Damit sind der Reihe nach die Koeffizienten
$\psi_{m-1},\psi_{m-2},\dots,\psi_2,\psi_1,\psi_0,A$
eindeutig festgelegt, was die Existenz des Polynoms $\psi$ und der
Konstanten $A$ beweist.
\end{proof}

Der Satz besagt, dass der erste Summand in~\eqref{buch:intalg:wurzel:eqn:summe}
zu einem Integranden führt, dessen Integral sich auf das Integral von
$1/\vartheta$ reduzieren lässt.
Diese Art von Integral tritt auch bei der Berechnung des zweiten Summanden
auf und wird im
Abschnitt~\ref{buch:intalg:wurzel:subsection:stammfunktiontheta}
berechnet.

%
% Partialbruchzeregung des gebrochenen Teils
%
\subsection{Partialbruchzerlegung des gebrochenen Teils}
Es bleibt jetzt also noch das Integral der Form
\[
\int \frac{w(z)}{\vartheta}\,dz
\]
zu berechnen.

%
% Die Stammfunktion von 1/theta
%
\subsection{Die Stammfunktion von $1/\vartheta$
\label{buch:intalg:wurzel:subsection:stammfunktiontheta}}
Durch Nachrechnen der Ableitung von
\[
\Theta
=
\frac{1}{\!\sqrt{a}}
\log\Bigl(
2az + b - 2\!\sqrt{a} \vartheta
\Bigr)
\]
bekommt man
\begin{align*}
\Theta'
&=
-\frac{1}{\!\sqrt{a}}
\cdot
\frac{1}{2az+b-2\!\sqrt{a}\vartheta}
\cdot
(2a - 2\!\sqrt{a}\vartheta')
\\
&=
-\frac{1}{\!\sqrt{a}}
\cdot
\frac{1}{2az+b-2\!\sqrt{a}\vartheta}
\cdot
\biggl(
2a -2\!\sqrt{a}\frac{2az+b}{2\vartheta}
\biggr)
\\
&=
-\frac{1}{\!\sqrt{a}}
\cdot
\frac{1}{2az+b-2\!\sqrt{a}\vartheta}
\cdot
\frac{
4a\vartheta -2\!\sqrt{a}(2az+b)
}{
2\vartheta
}
\\
&=
-\frac{1}{\!\sqrt{a}}
\cdot
\frac{1}{2az+b-2\!\sqrt{a}\vartheta}
\cdot
\frac{
2\!\sqrt{a}\cdot
\bigl(
\!\sqrt{a}\vartheta
-
2az-b
\bigr)
}{
2\vartheta
}
=
\frac{1}{\vartheta},
\end{align*}
Die Funktion $\Theta$ ist also eine Stammfunktion von $1/\vartheta$.
Sie liegt in einem Erweiterungskörper, dem die Konstante $\!\sqrt{a}$
und der Logarithmus $\log(2az+b-2\!\sqrt{a}\vartheta)$ hinzugefügt
werden ist.
Damit der Logarithmus hinzugefügt werden kann, muss auch noch $b$
hinzugefügt werden.
Somit ist
\[
\Theta
\in
\mathbb{Q}\bigl(z,\theta,\!\sqrt{a},b,\log(\!\sqrt{a}\vartheta-2az-b)\bigr)
\]
eine elementare Stammfunktion von $1/\vartheta$.


%
% 2-liouville.tex
%
% (c) 2026 Prof Dr Andreas Müller
%
\section{Das Prinzip von Liouville
\label{buch:intalg:section:liouville}}
\kopfrechts{Das Prinzip von Liouville}%
Die in Kapitel~\ref{chapter:ratint} durchgerechneten Beispiele 
uggerieren, dass eine Stammfunktion immer nur eine Linearkombination
von Logarithmen der Teilnenner der Partialbruchzerlegung ist.
Das Prinzip von Liouville zeigt, dass dies ganz allgemein gilt.

%
% Das Prinzip
%
\subsection{Das Prinzip 
\label{buch:intalg:liouville:subsection:prinzip}}

\begin{satz}
Sei $F$ ein Funktionenkörper mit Konstantenkörper $K$ und $f\in F$.
Wenn $f$ eine Stammfunktion $g$ in einer elementaren Erweiterung
$G\supset F$ hat, die den gleichen Konstantenkörper wie $F$ hat, dann
gibt es $v_0,v_1,\dots,v_n\in F$ und $c_1,\dots,c_k\in K$ mit
\begin{equation}
f
=
Dv_0
+
\sum_{i=1}^n c_i\frac{Dv_i}{v_i},
\end{equation}
d.~h.~
\[
\int f(z)\,dz
=
v_0
+ 
\sum_{i=1}^n c_i \log(v_i).
\]
$G$ muss also alle Logarithmen $\log(v_i)$, $i=1,\dots,n$, enthalten.
\end{satz}

Wie kann man überhaupt eine solche Aussage beweisen?
Wenn $g$ die Stammfunktion von $f$ ist, dann muss die Ableitung
von $g$ wieder $f$ ergeben, also $Dg=f$.
Wenn also $g$ in einer Erweiterung $F(\vartheta_1,\dots,\vartheta_n)$
von $F$ liegt, dann kann man versuchen, einen Ausdruck für $g$ in den
Elementen $\vartheta_k$ abzuleiten.
Dabei muss ein Ausdruck entstehen, welcher mit $f$ übereinstimmt.
Alle Terme mit den Elementen $\vartheta_k$ müssen sich daher wegheben.

Der allgemeine Fall einer Erweiterung der Form
$F(\vartheta_1,\dots,\vartheta_n)$ führt zu ziemlich komplizierten
Ausdrücken.
Daher beginnen wir damit einer Erweiterung $F(\vartheta)$ um ein
einziges $\vartheta$.
In diesem Fall muss $g$ die Form einer rationalen Funktion
\[
g
=
\frac{a(\vartheta)}{b(\vartheta)} 
=
\frac{
a_0 + a_1\vartheta + \dots a_n\vartheta^n
}{
b_0 + b_1\vartheta + \dots b_m\vartheta^m
},
\]
wobei $a(\vartheta)\in F[\vartheta]$ und $b(\vartheta)\in F[\vartheta]$
Polynome mit Koeffizienten in $F$ sind.
In diesem Fall lässt sich die Ableitung von $g$ leichter 
berechnen.
Zu diesem Zweck berechnen wir im
Abschnitt~\ref{buch:intalt:liouville:subsection:polynome}
zunächst die Ableitung von Polynomen von $\vartheta$ für die
verschiedenen Arten von Erweiterungen.
Diese Resultate verwenden wir dann im
Abschnitt~\ref{buch:intalt:liouville:subsection:spezialfaelle},
um das Prinzip von Liouville für eine Erweiterung der Form $F(\vartheta)$
zu beweisen.
In Abschnitt~\ref{buch:intalt:liouville:subsection:allgemein}
wird dann der allgemeine Fall bewiesen.

%
% Polynome von Erweiterungen
%
\subsection{Polynome von Erweiterungen
\label{buch:intalg:liouville:subsection:polynome}}

%
% Beweis des Prinzips von Liouville in Spezialfällen
%
\subsection{Beweis des Prinzips von Liouville in Spezialfällen
\label{buch:intalg:liouville:subsection:spezialfaelle}}

%
% Beweis im allgemeinen Fall
%
\subsection{Beweis im allgemeinen Fall
\label{buch:intalg:liouville:subsection:allgemein}}




\input{chapters/070-intalg/4-risch.tex}


%\uebungsabschnitt
%
%\aufgabetoplevel{chapters/070-integration/uebungsaufgaben}
%\begin{uebungsaufgaben}
%\uebungsaufgabe{701}
%\uebungsaufgabe{702}
%\end{uebungsaufgaben}
%\enduebungsabschnitt

